% Generated by manuscript/tables/sequence_redundancy/extract_redundancy.py
\begingroup
\small
\setlength{\tabcolsep}{3pt}
\renewcommand{\arraystretch}{1.15}
\setlength{\LTleft}{0pt}
\setlength{\LTright}{0pt}
\setlength{\LTcapwidth}{\linewidth}
\begin{longtable}{@{}>{\raggedright\arraybackslash}p{\dimexpr0.34\linewidth\relax}>{\centering\arraybackslash}p{\dimexpr0.11\linewidth-2\tabcolsep\relax}>{\centering\arraybackslash}p{\dimexpr0.1\linewidth-2\tabcolsep\relax}>{\centering\arraybackslash}p{\dimexpr0.15\linewidth-2\tabcolsep\relax}>{\centering\arraybackslash}p{\dimexpr0.17\linewidth-2\tabcolsep\relax}>{\centering\arraybackslash}p{\dimexpr0.13\linewidth-2\tabcolsep\relax}@{}}
\caption{\textbf{Sequence redundancy between the supplied training and test partitions.} Exact matches are full-length sequence equality after uppercasing and nucleotide U-to-T normalization. Similarity includes exact matches and MMseqs2 matches with at least 90\% identity and 80\% coverage of each sequence; nucleotide searches consider both strands. Counts refer to test records with at least one training match; percentages use the test-set size. Retained denotes records without a detected training match, not a newly constructed test set. A dash indicates no supplied test set. All 471 test records for Model 58 match training exactly. Search details and limitations are given in the Supplementary Notes (Sequence redundancy).}\label{tab:sequence_overlap}\\
\toprule
\textbf{Model} & \shortstack{\textbf{Training}\\\textbf{records}} & \shortstack{\textbf{Test}\\\textbf{records}} & \shortstack{\textbf{Exact}\\$n$ (\%)} & \shortstack{\textbf{Similar}\\\textbf{(incl. exact)}\\$n$ (\%)} & \shortstack{\textbf{Retained}\\\textbf{test records}} \\
\midrule
\endfirsthead
\toprule
\textbf{Model} & \shortstack{\textbf{Training}\\\textbf{records}} & \shortstack{\textbf{Test}\\\textbf{records}} & \shortstack{\textbf{Exact}\\$n$ (\%)} & \shortstack{\textbf{Similar}\\\textbf{(incl. exact)}\\$n$ (\%)} & \shortstack{\textbf{Retained}\\\textbf{test records}} \\
\midrule
\endhead
\midrule
\multicolumn{6}{r}{\textit{Continued on next page}}\\
\endfoot
\bottomrule
\endlastfoot
Model 1: antibody peptides & 63,400 & 471 & 0 (0.00) & 332 (70.49) & 139 \\
Model 2: anticancer peptides & 500 & 164 & 0 (0.00) & 8 (4.88) & 156 \\
Model 3: anticancer peptides & 344 & -- & -- & -- & -- \\
Model 4: anticancer peptides & 1,552 & 388 & 0 (0.00) & 85 (21.91) & 303 \\
Model 5: anticancer peptides & 1,378 & 344 & 0 (0.00) & 133 (38.66) & 211 \\
Model 6: anti-coronavirus & 465 & -- & -- & -- & -- \\
Model 7: anti-coronavirus and random non-secretory proteins & 648 & -- & -- & -- & -- \\
Model 8: antifungal peptides & 1,556 & 430 & 0 (0.00) & 50 (11.63) & 380 \\
Model 9: anti-hypertensive peptides & 1,826 & 772 & 244 (31.61) & 267 (34.59) & 505 \\
Model 10: antimalarial peptides & 653 & 163 & 0 (0.00) & 5 (3.07) & 158 \\
Model 11: antimalarial peptides & 1,819 & 455 & 0 (0.00) & 16 (3.52) & 439 \\
Model 12: antimicrobial peptides & 18,114 & 7,764 & 0 (0.00) & 9 (0.12) & 7,755 \\
Model 13: antimicrobial peptides & 3,284 & 1,840 & 0 (0.00) & 288 (15.65) & 1,552 \\
Model 14: antimicrobial peptides & 13,428 & 8,953 & 0 (0.00) & 849 (9.48) & 8,104 \\
Model 15: anti-MRSA strains peptides & 796 & 199 & 0 (0.00) & 0 (0.00) & 199 \\
Model 16: antioxidant proteins & 1,780 & 466 & 38 (8.15) & 40 (8.58) & 426 \\
Model 17: anti-parasitic peptides & 2,118 & 92 & 19 (20.65) & 19 (20.65) & 73 \\
Model 18: antiviral & 899 & 98 & 0 (0.00) & 49 (50.00) & 49 \\
Model 19: antiviral and random non-secretory proteins & 1,059 & 115 & 0 (0.00) & 40 (34.78) & 75 \\
Model 20: antiviral peptides & 4,642 & 1,246 & 0 (0.00) & 289 (23.19) & 957 \\
Model 21: bitter peptides & 512 & 128 & 0 (0.00) & 36 (28.12) & 92 \\
Model 22: blood-brain barrier peptides & 200 & 38 & 0 (0.00) & 0 (0.00) & 38 \\
Model 23: DNA-binding proteins & 1,075 & 183 & 75 (40.98) & 79 (43.17) & 104 \\
Model 24: DNA-binding proteins & 396 & -- & -- & -- & -- \\
Model 25: DNase I hypersensitive sites & 1,017 & -- & -- & -- & -- \\
Model 26: DPP IV inhibitory peptides & 1,064 & 266 & 3 (1.13) & 35 (13.16) & 231 \\
Model 27: hemolytic peptides & 884 & 220 & 0 (0.00) & 56 (25.45) & 164 \\
Model 28: recombination spots & 1,050 & -- & -- & -- & -- \\
Model 29: lncRNA (\textit{Homo sapiens}) & 16,000 & 5,000 & 21 (0.42) & 344 (6.88) & 4,656 \\
Model 30: lncRNA (\textit{Mus musculus}) & 8,400 & 3,600 & 28 (0.78) & 207 (5.75) & 3,393 \\
Model 31: lncRNA (\textit{Triticum aestivum}) & 8,000 & 4,000 & 8 (0.20) & 339 (8.47) & 3,661 \\
Model 32: lncRNA (\textit{Zea mays}) & 36,000 & -- & -- & -- & -- \\
Model 33: m5C sites (\textit{Arabidopsis thaliana}) & 10,578 & 2,000 & 0 (0.00) & 0 (0.00) & 2,000 \\
Model 34: m5C sites (\textit{Homo sapiens}) & 240 & -- & -- & -- & -- \\
Model 35: m5C sites (\textit{Mus musculus}) & 194 & -- & -- & -- & -- \\
Model 36: m5C sites (\textit{Saccharomyces cerevisiae}) & 422 & -- & -- & -- & -- \\
Model 37: neuropeptides & 3,880 & 970 & 0 (0.00) & 102 (10.52) & 868 \\
Model 38: non-classical secreted proteins & 587 & 68 & 0 (0.00) & 0 (0.00) & 68 \\
Model 39: peptide toxicity & 3,284 & 580 & 0 (0.00) & 34 (5.86) & 546 \\
Model 40: phage virion proteins & 500 & 126 & 5 (3.97) & 5 (3.97) & 121 \\
Model 41: proinflammatory peptides & 1,705 & 290 & 0 (0.00) & 0 (0.00) & 290 \\
Model 42: protein lysine crotonylation sites & 32,418 & 8,169 & 0 (0.00) & 0 (0.00) & 8,169 \\
Model 43: quorum-sensing peptides & 400 & 40 & 0 (0.00) & 7 (17.50) & 33 \\
Model 44: real microRNA precursors & 3,224 & -- & -- & -- & -- \\
Model 45: ribosome binding site sequences & 275,849 & 27,654 & 0 (0.00) & 7,556 (27.32) & 20,098 \\
Model 46: sigma70 promoters & 2,141 & -- & -- & -- & -- \\
Model 47: small non-coding RNA and shuffled sequences & 364 & -- & -- & -- & -- \\
Model 48: toehold switch sequences & 91,534 & 168 & 0 (0.00) & 0 (0.00) & 168 \\
Model 49: Tumor T cell antigens & 788 & 197 & 4 (2.03) & 22 (11.17) & 175 \\
Model 50: umami peptides & 353 & 89 & 1 (1.12) & 17 (19.10) & 72 \\
Model 51: lncRNAs subcellular localization - training set with 4 classes & 643 & 396 & 7 (1.77) & 16 (4.04) & 380 \\
Model 52: lncRNAs subcellular localization - training set with 5 classes & 600 & 396 & 35 (8.84) & 71 (17.93) & 325 \\
Model 53: mRNA subcellular localization & 12,410 & 2,499 & 0 (0.00) & 70 (2.80) & 2,429 \\
Model 54: non-coding RNA - 4 classes & 1,566 & 663 & 2 (0.30) & 152 (22.93) & 511 \\
Model 55: non-coding RNA - 8 classes & 2,396 & 654 & 0 (0.00) & 85 (13.00) & 569 \\
Model 56: non-coding RNA - E. coli K12 & 205 & 51 & 0 (0.00) & 27 (52.94) & 24 \\
Model 57: non-coding RNA - Multiple bacterial phyla & 783 & 549 & 63 (11.48) & 184 (33.52) & 365 \\
Model 58: antibody peptides & 67,769 & 471 & 471 (100.00) & 471 (100.00) & 0 \\
Model 59: ribosome binding site sequences & 275,849 & 27,654 & 0 (0.00) & 7,556 (27.32) & 20,098 \\
Model 60: toehold switch sequences & 91,534 & 168 & 0 (0.00) & 0 (0.00) & 168 \\
Model 61: protein solubility & 3,148 & -- & -- & -- & -- \\
Model 62: enzyme optimal catalytic temperature & 2,626 & 291 & 4 (1.37) & 30 (10.31) & 261 \\
Model 63: protein melting temperature & 2,255 & 251 & 0 (0.00) & 5 (1.99) & 246 \\
Model 64: yeast protein half-life & 3,118 & -- & -- & -- & -- \\
\end{longtable}
\endgroup
